\chapter{Why Ada for Microcontrollers}

Microcontroller firmware lives in a hostile place: a few hundred kilobytes of
RAM, no operating system to catch mistakes, interrupts that fire between any two
instructions, and hardware registers whose exact bit layout is law. A bug is not
a stack trace on a screen. It is a board that silently locks up in the field.
Ada was designed for exactly this kind of software, and the language gives the
firmware author tools that C simply does not have.

\section{The case for Ada}\index{type safety}\index{representation clause}\index{memory-mapped I/O}\index{private child}\index{contracts}

\begin{description}[style=nextline]
\item[Strong, range-checked types.] In Ada a ``GPIO pin'' need not be an
  \texttt{int}. It can be a type that admits \emph{only} the pins the silicon
  actually exposes; passing a flash-mapped pad becomes a compile error, not a
  hung chip. Numeric types carry ranges, and the compiler (or a run-time check)
  enforces them.
\item[Representation control.] Ada lets you state, in the language, the precise
  bit and byte layout of a record and the exact address of an object. A hardware
  register is just an Ada record at a fixed address: no casts, no macros, no
  guessing about padding (Sections~\ref{ch:registers} and~\ref{ch:bitfields} of
  the Talking to Hardware chapter).
\item[Packages and strong interfaces.] Drivers are packages with a visible
  specification and a hidden body. The unsafe register pokes can be sealed inside
  a \emph{private child} that application code literally cannot name.
\item[Tasking and real-time in the language.] Tasks, protected objects (data
  guarded by a monitor), \texttt{delay until}, and priorities are part of Ada
  itself, not a bolted-on RTOS API. On this port they are implemented by a
  native Ada runtime that runs directly on the silicon.
\item[Restriction profiles.] With one pragma you can forbid features that do not
  belong in certified or tiny firmware (no heap, no recursion in the runtime, no
  unbounded constructs) and have the compiler \emph{prove} you stayed within the
  subset -- the Ravenscar and Jorvik profiles (Chapter~\ref{ch:profiles}).
\item[Contracts.] Preconditions, postconditions and type predicates let you
  state intent (``this key must be 128 or 256 bits'') and have it checked.
\end{description}

The result is firmware where a large class of mistakes (wrong-width register
writes, out-of-range pins, aliasing of a peripheral by two tasks, forgetting to
release a DMA channel) are caught by the compiler or by a cheap run-time
check, long before they reach silicon.

\section{What this book covers}

This book documents a complete bare-metal stack for the dual-core ESP32-S3
(Xtensa LX7) written almost entirely in Ada, running directly on the silicon
with its own runtime, boot path and tooling:

\begin{enumerate}
\item the Ada language techniques that make register and bit-field programming
  safe and pleasant (the Ada Language Techniques and Talking to Hardware
  chapters, Chapters~\ref{ch:embedded} and~\ref{ch:hardware});
\item the anatomy of an Ada application, from \texttt{Main} down through the
  runtime to the boot ROM, and the three selectable runtime profiles
  (Chapters~\ref{ch:anatomy}--\ref{ch:build});
\item the reusable peripheral HAL (twenty-plus task-safe drivers) with the
  API and worked examples for each (Chapters~\ref{ch:hal}--\ref{ch:hal3});
\item mass storage: the SD-card block drivers and the pure-Ada ext2/3/4
  filesystem (Chapter~\ref{ch:storage}).
\end{enumerate}

New to Ada? Chapter~\ref{ch:ada-primer}, \emph{Ada for C Programmers}, is a
one-sitting crash course in the slice of the language this book uses
(declarations, subprograms and parameter modes, packages, control flow), framed
against the C you already know; the rest of the book assumes that much. Individual
Ada terms that may still be unfamiliar to a reader coming from C
(\emph{library-level}, \emph{elaboration}, \emph{aspect}, \emph{protected object}
and the like) are also collected in the Terminology section of the Ada Language
Techniques chapter (\S\ref{sec:terms}).

